% ============================================================================
% Welcome Page — Grade 7 Math Step by Step
% Layered header with soft gradient-like fills, arrow/chevron motifs,
% centred welcome prose, and a 2×2 showcard grid. THE TPT thumbnail page.
% UNIQUE to Step-by-Step: arrow/chevron motifs, orange brand colour.
% ============================================================================

\thispagestyle{empty}

% ── TikZ full-page background ────────────────────────────────────────────
\begin{tikzpicture}[remember picture, overlay]
    % Soft warm wash
    \fill[funOrange!2]
        (current page.south west) rectangle (current page.north east);

    % ── Layered header — three soft overlapping bands ──
    % Outermost: wide orange band with angled bottom
    \fill[funOrange!7]
        (current page.north west) --
        (current page.north east) --
        ([yshift=-9.8cm]current page.north east) --
        ([yshift=-10.6cm]current page.north west) -- cycle;

    % Middle: slightly inset teal tint for depth
    \fill[funTeal!5]
        ([xshift=0.6cm, yshift=-0.6cm]current page.north west) --
        ([xshift=-0.6cm, yshift=-0.6cm]current page.north east) --
        ([xshift=-0.6cm, yshift=-9.4cm]current page.north east) --
        ([xshift=0.6cm, yshift=-10.1cm]current page.north west) -- cycle;

    % Inner: lighter lift
    \fill[white, opacity=0.35]
        ([xshift=1.6cm, yshift=-1.6cm]current page.north west) --
        ([xshift=-1.6cm, yshift=-1.6cm]current page.north east) --
        ([xshift=-1.6cm, yshift=-8.8cm]current page.north east) --
        ([xshift=1.6cm, yshift=-9.4cm]current page.north west) -- cycle;

    % ── Chevron cluster — top left ──
    \node[isosceles triangle, isosceles triangle apex angle=60,
          fill=funOrange!12, minimum size=8mm, inner sep=0pt, rotate=-30]
        at ([xshift=2.4cm, yshift=-1.6cm]current page.north west) {};
    \node[isosceles triangle, isosceles triangle apex angle=60,
          fill=funTeal!8, minimum size=5mm, inner sep=0pt, rotate=-30]
        at ([xshift=3.8cm, yshift=-2.7cm]current page.north west) {};
    \node[isosceles triangle, isosceles triangle apex angle=60,
          fill=funOrange!7, minimum size=3.5mm, inner sep=0pt, rotate=-30]
        at ([xshift=5.4cm, yshift=-1.3cm]current page.north west) {};

    % ── Chevron cluster — top right ──
    \node[isosceles triangle, isosceles triangle apex angle=60,
          fill=funTeal!9, minimum size=6.5mm, inner sep=0pt, rotate=-30]
        at ([xshift=-2.6cm, yshift=-1.9cm]current page.north east) {};
    \node[isosceles triangle, isosceles triangle apex angle=60,
          fill=funOrange!10, minimum size=4mm, inner sep=0pt, rotate=-30]
        at ([xshift=-4.4cm, yshift=-2.9cm]current page.north east) {};
    \node[isosceles triangle, isosceles triangle apex angle=60,
          fill=funOrange!6, minimum size=2.8mm, inner sep=0pt, rotate=-30]
        at ([xshift=-6.2cm, yshift=-1.5cm]current page.north east) {};

    % ── Chevron cluster — bottom corners ──
    \node[isosceles triangle, isosceles triangle apex angle=60,
          fill=funTeal!5, minimum size=5mm, inner sep=0pt, rotate=-30]
        at ([xshift=2.6cm, yshift=2.3cm]current page.south west) {};
    \node[isosceles triangle, isosceles triangle apex angle=60,
          fill=funOrange!5, minimum size=3.5mm, inner sep=0pt, rotate=-30]
        at ([xshift=5.2cm, yshift=1.7cm]current page.south west) {};
    \node[isosceles triangle, isosceles triangle apex angle=60,
          fill=funOrange!6, minimum size=4.5mm, inner sep=0pt, rotate=-30]
        at ([xshift=-3cm, yshift=2.6cm]current page.south east) {};
    \node[isosceles triangle, isosceles triangle apex angle=60,
          fill=funTeal!4, minimum size=2.5mm, inner sep=0pt, rotate=-30]
        at ([xshift=-5.4cm, yshift=1.5cm]current page.south east) {};

    % ── Tracking text ──
    \node[font=\fontsize{9}{11}\selectfont\sffamily\color{funOrangeDark!50}]
        at ([yshift=-2.4cm]current page.north)
        {\textls[180]{YOUR STEP-BY-STEP MATH GUIDE}};

    % ── Main title ──
    \node[font=\fontsize{34}{40}\selectfont\sffamily\bfseries\color{black!88},
          text width=0.82\paperwidth, align=center]
        at ([yshift=-4.6cm]current page.north)
        {Math Step by Step};

    % ── Subtitle ──
    \node[font=\normalsize\sffamily\color{black!45},
          text width=0.7\paperwidth, align=center]
        at ([yshift=-6.2cm]current page.north)
        {Clear Steps \;\textbullet\; Worked Examples \;\textbullet\;
        Guided Practice};

    % ── Elegant separator: triple-arrow motif ──
    \draw[funOrange!22, line width=0.9pt]
        ([xshift=3.5cm, yshift=-7.5cm]current page.north west) --
        ([xshift=-3.5cm, yshift=-7.5cm]current page.north east);
    % Central chevron
    \node[isosceles triangle, isosceles triangle apex angle=60,
          fill=funOrange!15, draw=funOrange!28, line width=0.4pt,
          minimum size=5.5mm, inner sep=0pt, rotate=-30]
        at ([yshift=-7.5cm]current page.north) {};
    % Flanking chevrons
    \node[isosceles triangle, isosceles triangle apex angle=60,
          fill=funTeal!8, minimum size=3.5mm, inner sep=0pt, rotate=-30]
        at ([xshift=-1.2cm, yshift=-7.5cm]current page.north) {};
    \node[isosceles triangle, isosceles triangle apex angle=60,
          fill=funTeal!8, minimum size=3.5mm, inner sep=0pt, rotate=-30]
        at ([xshift=1.2cm, yshift=-7.5cm]current page.north) {};

    % ── Thin outer frame ──
    \draw[black!6, line width=0.3pt]
        ([xshift=8mm, yshift=-8mm]current page.north west) rectangle
        ([xshift=-8mm, yshift=8mm]current page.south east);

    % ── Copyright footer ──
    \node[font=\sffamily\tiny\color{black!32}, anchor=south]
        at ([yshift=0.5cm]current page.south)
        {\textcopyright~2026 Dr.\ A.\ Nazari \;\textbullet\; ViewMath.com};

\end{tikzpicture}

% ── Welcome Text ──────────────────────────────────────────────────────────
\vspace*{7.6cm}

\begin{center}
\begin{minipage}{0.76\textwidth}
    \centering\sffamily\large\color{black!68}
    \setlength{\parskip}{6pt}

    This book breaks every Grade~7 math topic into clear,
    numbered steps. Read the steps, follow the worked examples,
    try the guided practice, then solve problems on your own.
    The answer key in the back lets you check your work
    and learn from every mistake.
\end{minipage}
\end{center}

% ── Feature Showcards — 2×2 grid ─────────────────────────────────────────
\vspace{7mm}

\begin{center}
\begin{minipage}{0.86\textwidth}

\noindent
\begin{minipage}[t]{0.48\textwidth}
\stShowcard{funOrange}{\faRoute}{Step-by-Step Method}{%
    Every topic is taught as a numbered procedure ---
    follow the steps and you'll get the answer.}
\end{minipage}%
\hfill
\begin{minipage}[t]{0.48\textwidth}
\stShowcard{funTeal}{\faBookOpen}{9 Chapters, 56 Topics}{%
    From ratios and percents to probability ---
    every Grade~7 standard covered.}
\end{minipage}

\vspace{1mm}

\noindent
\begin{minipage}[t]{0.48\textwidth}
\stShowcard{funBlue}{\faIcon{pencil-alt}}{Guided Practice}{%
    Try-it-yourself problems with blanks that walk you
    through each step before you go solo.}
\end{minipage}%
\hfill
\begin{minipage}[t]{0.48\textwidth}
\stShowcard{funGreen}{\faCheckDouble}{Complete Answer Key}{%
    Every problem has an answer in the back. Check your work,
    learn from mistakes, and try again.}
\end{minipage}

\end{minipage}
\end{center}

\clearpage
